% from ega3-1-fr.tex
\oldpage[III]{147}

entraîne a), on peut se borner encore au cas où
$\mathfrak Y=\operatorname{Spf}(B)$, $\mathfrak X=\operatorname{Spf}(A)$,
$A$ étant une $B$-algèbre finie~; comme $A/\mathfrak RA$ est alors une
$B/\mathfrak R$-algèbre finie, il résulte de (I, 10.10.9) que les conditions
de a) sont satisfaites.
\end{proof}

\begin{definition}[4.8.2]
\label{III.4.8.2-fr}
Lorsque les propriétés équivalentes a), b), c) de
(\hyperref[III.4.8.1-fr]{4.8.1}) sont vérifiées, on dit que le morphisme $f$
est \emph{fini}, ou que $\mathfrak X$ est un \emph{$\mathfrak Y$-préschéma
formel fini}, ou un préschéma formel \emph{fini au-dessus de $\mathfrak Y$}.
\end{definition}

\begin{proposition}[4.8.3]
\label{III.4.8.3-fr}
\begin{enumerate}
\item[(i)] Une immersion fermée de préschémas formels localement noethériens
  est un morphisme fini.
\item[(ii)] Le composé de deux morphismes finis de préschémas formels
  localement noethériens est un morphisme fini.
\item[(iii)] Soient $\mathfrak X$, $\mathfrak Y$, $\mathfrak S$ trois
  préschémas formels localement noethériens, $f:\mathfrak X\to\mathfrak S$ un
  morphisme fini, $g:\mathfrak Y\to\mathfrak S$ un morphisme~; alors le
  morphisme
  $\mathfrak X\times_{\mathfrak S}\mathfrak Y\to\mathfrak Y$ est fini.
\item[(iv)] Soient $\mathfrak S$ un préschéma formel localement noethérien,
  $\mathfrak X'$, $\mathfrak Y'$ deux préschémas formels localement
  noethériens tels que $\mathfrak X'\times_{\mathfrak S}\mathfrak Y'$ soit
  localement noethérien. Si $\mathfrak X$, $\mathfrak Y$ sont des
  $\mathfrak S$-préschémas formels localement noethériens,
  $f:\mathfrak X\to\mathfrak X'$, $g:\mathfrak Y\to\mathfrak Y'$ deux
  $\mathfrak S$-morphismes finis, alors $f\times_{\mathfrak S}g$ est un
  morphisme fini.
\item[(v)] Soient $f:\mathfrak X\to\mathfrak Y$,
  $g:\mathfrak Y\to\mathfrak Z$ deux morphismes de préschémas formels
  localement noethériens tels que $g$ soit de type fini et séparé~; alors, si
  $g\circ f$ est un morphisme fini, $f$ est un morphisme fini.
\end{enumerate}
\end{proposition}

\begin{proof}
(i) est trivial, et les autres assertions se ramènent aussitôt aux
propositions correspondantes pour les morphismes de préschémas usuels
(II, 6.1.5) à l'aide du critère a) de
(\hyperref[III.4.8.1-fr]{4.8.1})~; nous laissons les détails au lecteur, sur le
modèle de (I, 10.13.5).
\end{proof}

\begin{corollary}[4.8.4]
\label{III.4.8.4-fr}
Sous les hypothèses de (I, 10.9.9), si $f$ est un morphisme fini, il en est de
même de son prolongement $\widehat f$ aux complétés.
\end{corollary}

\begin{corollary}[4.8.5]
\label{III.4.8.5-fr}
Si $\mathfrak X$ est un préschéma formel fini au-dessus de $\mathfrak Y$,
$f:\mathfrak X\to\mathfrak Y$ le morphisme structural, alors, pour tout ouvert
$U\subset\mathfrak Y$, $f^{-1}(U)$ est fini au-dessus de $U$.
\end{corollary}

\begin{proposition}[4.8.6]
\label{III.4.8.6-fr}
Si $f:\mathfrak X\to\mathfrak Y$ est un morphisme fini de préschémas formels
localement noethériens, $f_*(\mathcal O_{\mathfrak X})$ est une
$\mathcal O_{\mathfrak Y}$-Algèbre cohérente.
\end{proposition}

\begin{proof}
On peut considérer $f$ comme limite inductive d'un système inductif $(f_n)$ de
morphismes $f_n:X_n\to Y_n$~; montrons que les $f_n$ sont des morphismes finis
et que $f_*(\mathcal O_{\mathfrak X})$ est isomorphe à la limite projective des
$(f_n)_*(\mathcal O_{X_n})$, ce qui établira notre assertion (I, 10.10.5). Il
suffit de se borner au cas où $\mathfrak Y=\operatorname{Spf}(B)$,
$\mathfrak X=\operatorname{Spf}(A)$, et de remarquer que si $\mathfrak R$ est
un idéal de définition de $B$ et $A$ un $B$-module de type fini,
$A/\mathfrak R^{n+1}A$ est un module de type fini sur
$B/\mathfrak R^{n+1}B$, et que $A$ est limite projective des
$A/\mathfrak R^{n+1}A$.
\end{proof}

Réciproquement~:

\begin{proposition}[4.8.7]
\label{III.4.8.7-fr}
Soient $\mathfrak Y$ un préschéma formel localement noethérien, $\mathcal A$
une $\mathcal O_{\mathfrak Y}$-Algèbre cohérente. Il existe un préschéma
formel $\mathfrak X$ fini au-dessus de $\mathfrak Y$, défini à un
$\mathfrak Y$-isomorphisme unique près, et tel que
$f_*(\mathcal O_{\mathfrak X})=\mathcal A$, $f:\mathfrak X\to\mathfrak Y$
étant le morphisme structural.
\end{proposition}

\begin{proof}
Soit $\mathcal K$ un Idéal de définition de $\mathfrak Y$, et posons
\[
Y_n=(\mathfrak Y,\mathcal O_{\mathfrak Y}/\mathcal K^{n+1}),\qquad
\mathcal A_n=\mathcal A/\mathcal K^{n+1}\mathcal A~;
\]
il est clair que $\mathcal A_n$ est une $\mathcal O_{Y_n}$-Algèbre finie et
définit donc un

